% Base: old/main.tex lines 1-3 (the initial manuscript's abstract and keywords).
% Final ladder: results/julia/final/ladder_3d_final.csv (trucks 12-156, complaints 0,
% time_limit 900-3600, gap_pct 0.0075 / 0.1628 / 0.022 / 14.5186 / 10.686); minimum
% fleets on rungs 1-3 (5, 10, 12 periods): requirement prunes, ladder_3d_final.log
% lines 15, 651, 705.  Volumetric stage: Table tab:volumetric (results/julia/
% parity_vol.csv against results/compact_ref/*.json): 68139.21 vs 68139.17
% (0.0001%), certified bound 0.0113% below, four memory-stop rows within 0.31%.
% Benchmark: results/julia/final/bench36/bench36_joined.csv (branch final-3d,
% 360c2ea): 32 rows with method_cost, median gap_pct 6.4642 against best_lb;
% bench36_method.csv complaints 0 and pack_undecided 0 on every row.
% Volumetric plans: Table tab:repack and census_compact36.json (unplaceable
% loads on 35 of 36 plans).
% Catalogue: results/catalogue_realism.json (kappa: nine at 1.000, two at 1.176,
% four at 1.231-2.000, so six of fifteen above 1).
% Words: 239 from the PDF with "3 600" as one word (240 as two).
We propose a mixed integer linear programming model, 3D-SCPTOP, for procurement transport planning, integrating procurement, inventory and transport decisions of a plant and its supplier with three-dimensional full-truckload constraints and closed-loop management of returnable transport items. We reformulate it by Dantzig--Wolfe decomposition, with the loaded truck as the unit of decision: every geometric, stacking, crush, payload and full-truckload constraint lives inside one pricing problem per direction, and a column keeps its layout, so plans are executable by construction. The master is solved by branch and price over the fleet vector with a flow-first relax-and-fix construction as primal heuristic, and a certified volumetric bound is a valid lower bound for the 3D problem. On the volumetric problem the method returns a plan $0.0001\%$ above a proved optimum. On the 3D problem, where the compact model returns no incumbent beyond five periods, it returns verified plans on all five instances of a ladder reaching twenty-four products and twenty periods, from $12$ to $156$ trucks within $900$ to $3\,600$ s, with provably minimum fleets through twelve periods and gaps to its certified bound of $0.0075\%$ to $0.163\%$, and $14.5\%$ and $10.7\%$ where a per-stack crush limit binds. On a 36-instance benchmark it returns $32$ verified plans, at a median gap of $6.46\%$ to the best valid bound. Volumetric plans prescribe loads proved impossible to place, so a volumetric capacity is not a conservative simplification: it over-credits six of fifteen industrial container types.

\keywords{ Procurement, inventory, transport planning,  full-truckload, 3D container loading }
